\section*{Chapter \ref{ch:rc:structural-credit-models} --- Structural Credit Models}

\begin{solution}{exo:rc:structural-credit-models:1}
Equity 18.89, debt $50-18.89 = 31.11$, a yield spread of 103 basis points
($-\ln(31.11/40)/5-0.04$), and a risk-neutral \omterm{def:rc:reduced-form-credit:pd}{probability of default} $N(-d_2) = 23.5\%$.
\end{solution}

\begin{solution}{exo:rc:structural-credit-models:2}
Assets follow a diffusion: to default within a few months they must fall from well above
the debt to below it, which a continuous path with a volatility of 10--20\% a year almost
never does. The probability, and the spread, vanish faster than linearly as the maturity
shrinks.
\end{solution}

\begin{solution}{exo:rc:structural-credit-models:3}
$\mathrm{DD} = (\ln(40.69/40)+(0.08-\tfrac12\times0.156^2))/0.156 = 0.54$, a \omterm{def:rc:reduced-form-credit:pd}{probability of
default} of 29.3\%, higher than the five-year 15.4\%: over one year the drift has had little
time to lift the assets away from the face, and treating the whole face as due in one year
overstates the risk. Practical distance-to-default measures use a default point of
short-term debt plus part of the long-term debt for this reason.
\end{solution}

\begin{solution}{exo:rc:structural-credit-models:4}
With $\nu = 0$ the formula reduces to the reflection principle: $2N(\ln1.5/(0.3\sqrt2))-1 =
66.1\%$.
\end{solution}

\begin{solution}{exo:rc:structural-credit-models:5}
The hedge is linear in the notional: twice USD~1.74 million, a short of USD~3.49 million
of stock.
\end{solution}

\begin{solution}{exo:rc:structural-credit-models:6}
$\sigma_E = N(d_1)\sigma_VV/E$: equity is a levered claim on the assets, so a given
percentage move in the assets is a larger percentage move in the equity, by the ratio
$N(d_1)V/E$ (3.2 for the chapter's firm). When equity rises, leverage falls and the ratio
falls: a fixed asset volatility means an equity volatility that falls as the price rises,
the leverage effect.
\end{solution}

\begin{solution}{exo:rc:structural-credit-models:7}
\texttt{implied\_barrier()} gives 78.7\% of the face; the ten-year model spread with that
barrier is 380 basis points, below the five-year 500 because the curve is humped.
\end{solution}

\begin{solution}{exo:rc:structural-credit-models:8}
The \omterm{def:rc:structural-credit-models:merton}{Merton model} is known to give spreads far below the market's for most firms (no
short-dated default, no jumps, no liquidity or risk premia beyond the diffusion's), so a
large gap is the norm, not a signal. A capital-structure trade compares the market with a
model calibrated to the firm's own history (the barrier), and trades changes in the gap,
hedged; and a trade based on the level still needs the hedge and a view on why the gap
should close.
\end{solution}

\begin{solution}{pb:rc:structural-credit-models:1}
\textbf{1.} 229 and 500 basis points: the default swap is 271 basis points wide of the
model, a signal to sell protection.
\textbf{2.} $\partial s/\partial E = -0.0047$ per USD billion; short USD~1.74 million of
stock.
\textbf{3.} The 500-basis-point coupon on USD~10 million (USD~500\,000 a year, before the
cost of the equity short and its dividends, less funding income on the short's proceeds).
\textbf{4.} That the firm's credit is cheap relative to its equity: either spreads tighten
or equity falls.
\textbf{5.} Default (jump-to-default of USD~6 million on the protection against a stock
that falls to near zero), the model's convexity, volatility, idiosyncratic moves in one
leg (a takeover, a downgrade), and liquidity.
\textbf{6.} 161 basis points.
\textbf{7.} $500+(161-229) = 432$ basis points.
\textbf{8.} $-0.18\times1.74 \text{ million} = -\text{USD}~313\,961$.
\textbf{9.} $+\text{USD}~259\,002$ on the default swap, a hedged P\&L of $-\text{USD}~54\,958$
(the convexity of \cref{fig:rc:structural-credit-models:equityspread}).
\textbf{10.} $-\text{USD}~689\,903$ on the default swap; a total of $-\text{USD}~1\,003\,863$.
\textbf{11.} $-0.0030$ per USD billion: the slope flattens as equity rises, so the hedge
needs to be cut.
\textbf{12.} The asset volatility rises to 17.3\% and the model spread goes to 251 basis
points: in that version the model predicted a wider, not tighter, spread after the rally.
\textbf{13.} A takeover financed with debt, or a buyer of the shares whose interests differ
from the creditors', raises equity value and credit risk together: value moves from
creditors to shareholders, which a model with fixed debt and a single asset value cannot
produce.
\textbf{14.} 418 basis points: the default swap would have needed to tighten by 82 basis
points to pay for the equity leg's loss.
\textbf{15.} Nothing in the capital structure hedges a downgrade driven by the rating
itself; index protection or protection on similar names hedges part of it, and a smaller
position limits it.
\textbf{16.} Many held similar positions sized by similar models, so their unwinds moved
prices against each other, as the BIS review described.
\textbf{17.} By the stress loss of both legs moving against the position together, not by the
model's hedged risk; with limits per name and on event risk.
\textbf{18.} Close if the event changed the firm (the debt will be downgraded further, the
investor intends to change the capital structure); keep if the gap reflects forced
selling and the firm is unchanged, with funding to wait.
\textbf{19.} Named result: \emph{the capital-structure trade in May 2005}: the equity leg
loses USD~313\,961 and the default swap USD~689\,903, a total of USD~1\,003\,863 on
USD~10 million of protection over two days.
\textbf{20.} A link between the equity and credit prices of the same firm, for hedges,
signals and names without liquid default swaps.
\end{solution}

\begin{solution}{iq:rc:structural-credit-models:1}
Shareholders have limited liability: at the debt's maturity they repay and keep the rest,
or walk away and leave the assets to creditors. Their payoff is $\max(V_T-D,0)$, a call on
the assets struck at the debt. Debt is $\min(V_T,D)$: a riskless bond minus a put on the
assets.
\iqlookfor{limited liability, the call, and debt as bond minus put.}
\end{solution}

\begin{solution}{iq:rc:structural-credit-models:2}
Two equations in two unknowns: the equity value as a call on the assets, and the equity
volatility as $N(d_1)\sigma_VV/E$; solve by Newton. Alternatives iterate on the asset return
series implied by the daily equity values. The default point and horizon are choices, and
the equity volatility can come from options.
\iqlookfor{the two equations and the choices behind them.}
\end{solution}

\begin{solution}{iq:rc:structural-credit-models:3}
Continuous asset paths cannot reach the default point quickly, so short-dated default
probabilities are near zero. Fixes: jumps in the assets, an uncertain \omterm{def:rc:structural-credit-models:fp}{default barrier}
(the firm does not know exactly where it defaults), incomplete information about the
assets, or a hybrid with a hazard rate.
\iqlookfor{the diffusion reason and one or two fixes.}
\end{solution}

\begin{solution}{iq:rc:structural-credit-models:4}
The number of asset standard deviations between the expected asset value and the default
point at a horizon. Practitioners map it to default frequencies with an empirical table
from historical defaults, because the normal tail understates default rates, and use a
default point between short-term debt and total debt.
\iqlookfor{the formula and the empirical mapping.}
\end{solution}

\begin{solution}{iq:rc:structural-credit-models:5}
Sell protection and short equity when the default swap is wide of a \omterm{def:rc:structural-credit-models:structural}{structural model}
calibrated to the firm (or the reverse), with the equity amount set by the model's
spread sensitivity to equity times the position's \omterm{def:rc:reduced-form-credit:cs01}{CS01}. It loses on both legs if the
equity rallies (the short loses) while spreads widen (the protection loses), as when a
leveraged buyout or a takeover bid moves value from creditors to shareholders, or a
downgrade hits the debt alone.
\iqlookfor{the hedge ratio and the both-legs scenario.}
\end{solution}

\begin{solution}{iq:rc:structural-credit-models:6}
\omterm{def:rc:reduced-form-credit:reduced}{Reduced-form models} fit the credit curve exactly and price credit derivatives
consistently; they say nothing about why default happens. \omterm{def:rc:structural-credit-models:structural}{Structural models} explain
default from the balance sheet and link credit to equity, but fit market spreads poorly.
Use reduced-form for pricing and risk of credit instruments; structural for names
without default-swap quotes, equity-credit relative value, and credit scoring.
\iqlookfor{fit versus explanation, and the use of each.}
\end{solution}
